\section{Introduction}
A transformer of $L$ layers is conventionally an architectural commitment:
once trained, it is also evaluated at $L$ layers. The Neural Ordinary
Differential Equation (Neural ODE) framework~\cite{chen2018node} suggests
an alternative: a single block $f_\theta$ defines a vector field, and the
forward pass solves
\begin{equation}
\frac{dh}{dt}=f_\theta(h),\quad h(0)=\text{embed}(x),\quad h(1)=\text{logits}.
\end{equation}
The number of function evaluations (NFE) becomes a solver budget chosen at
inference time. We ask:
\begin{quote}\emph{Does a continuous-depth transformer obtain a more
favourable latency/accuracy trade-off than discrete-depth baselines, or
at least a more robust one?}\end{quote}

\paragraph{Findings.} On a synthetic copy task, both formulations reach
similar accuracy. The crucial difference is robustness to inference-time
budget: across seven solver settings the ODE next-token cross-entropy
varies by only $\OdeLossRange$ nats; discrete depth, when run at a depth
other than its training depth, varies by $\DiscreteLossRange$ nats.
Continuous depth therefore yields an \emph{anytime} model in the
algorithmic sense.

\paragraph{Contributions.} (i) A clean ODE-block transformer; (ii) a
solver-budget sweep covering Euler, RK4 and adaptive Dormand-Prince
schemes; (iii) a head-to-head comparison with discrete depths showing
that the loss-range gap is two orders of magnitude in favour of
continuous depth.
